% ----------------------------------------------------------------------
%  1.3. Симметрия кристаллов: 1.3.1–1.3.4 (оси симметрии, элементы теории
%  групп)
% ----------------------------------------------------------------------
\begin{orientir}
Какие повороты и отражения допускает кристалл и почему их так мало
(1.3.1--1.3.3); как их комбинации образуют группы (1.3.4); 32 точечные
группы и 14 решёток Браве (1.3.5); 230 пространственных групп с винтовыми
осями и плоскостями скольжения (1.3.6); предельные группы Кюри, нужные
для описания симметрии физических свойств (1.3.7).
\end{orientir}

Симметрия кристаллического пространства задаётся всеми
преобразованиями, которые сохраняют расстояния между точками и
совмещают пространство с самим собой. Элементы симметрии бывают:
\begin{itemize}
\item \term{открытые}~— содержат трансляцию (переносы, винтовые оси,
плоскости скольжения); они описывают симметрию бесконечного
пространства;
\item \term{закрытые}~— оставляют неподвижной хотя бы одну точку и после
конечного числа повторений возвращают пространство в исходное
положение. Они задаются ортогональными матрицами $R$ с $\det R=\pm1$ и
сводятся к \emph{поворотным осям} (собственные повороты, $\det R=+1$) и
\emph{инверсионным осям} (поворот с последующей инверсией в точке на оси,
$\det R=-1$).
\end{itemize}

\subsection{Поворотные оси симметрии}\label{ssec:rot-axes}

\begin{opredelenie}[Ось симметрии]
Прямая, при повороте вокруг которой на угол $\alpha_n=360^\circ/n$
кристаллическое пространство совмещается с собой. Целое $n$~—
\term{порядок оси}, $\alpha_n$~— её элементарный угол: через $n$ поворотов
пространство возвращается в исходное положение.
\end{opredelenie}

\paragraph{Почему возможны только оси 1, 2, 3, 4 и 6}
При повороте узел решётки $\vb x$ переходит в узел $\vb x'=R\vb x$. В
кристаллографической системе координат координаты узлов~— целые числа,
значит, и элементы матрицы $R$, и её след~— целые: $\operatorname{tr}R=N$. След не
зависит от системы координат (п.~\ref{ssec:orth}); в ортогональной системе
с осью $x$ вдоль оси поворота
\begin{equation}\label{eq:Rx}
R'=\begin{pmatrix}1&0&0\\0&\cos\alpha&-\sin\alpha\\0&\sin\alpha&\cos\alpha\end{pmatrix},\qquad
\operatorname{tr}R'=1+2\cos\alpha=N .
\end{equation}
Так как $|\cos\alpha|\le1$, $N$ может быть лишь $-1$, $0$, $1$, $2$, $3$
(табл.~\ref{tab:axes}).

\begin{table}[ht]
\centering
\caption{Возможные поворотные оси}\label{tab:axes}
\small
\begin{tabular}{@{}ccccc@{}}
\toprule
$\operatorname{tr}R$ & $-1$ & $0$ & $+1$ & $+2$ \ \ ($+3$) \\
\midrule
$\cos\alpha$ & $-1$ & $-1/2$ & $0$ & $1/2$ \ \ ($1$) \\
$\alpha$ & $180^\circ$ & $120^\circ$ & $90^\circ$ & $60^\circ$ \ \ ($0^\circ$) \\
Порядок оси $n$ & 2 & 3 & 4 & 6 \ \ (1) \\
\bottomrule
\end{tabular}
\end{table}

\begin{glavnoe}[Кристаллографическое ограничение]
В кристалле возможны только оси симметрии 1, 2, 3, 4 и 6-го порядков.
Осей 5-го, 7-го и более высоких порядков нет~— их запрещает
трансляционная симметрия (пятиугольниками нельзя без зазоров замостить
плоскость). Оси 5-го порядка появляются лишь в четырёхмерной
кристаллографии и у квазикристаллов (п.~\ref{ssec:space-groups}).
\end{glavnoe}

Действие осей удобно изображать на стереографической проекции: одно
направление общего положения и все направления, в которые его
переводят повороты (рис.~\ref{fig:axes-proj}, верхний ряд). В кубическом
кристалле оси 4-го порядка идут вдоль $\nfam{100}$, 3-го~— вдоль
$\nfam{111}$, 2-го~— вдоль $\nfam{110}$ (рис.~\ref{fig:cube-axes}).

\begin{figure}[ht]
\centering
\begin{tikzpicture}[scale=0.82]
  % создано tools/gen_stereo.py: действие осей симметрии
\draw[atomLine, line width=0.6pt] (0.0,0) circle (1);
\fill[fmBlue] (0.583,0.212) circle (0.075);
\node[font=\small] at (0.0,-1.32) {$1$};
\draw[atomLine, line width=0.6pt] (2.6,0) circle (1);
\fill[fmBlue] (3.183,0.212) circle (0.075);
\fill[fmBlue] (2.017,-0.212) circle (0.075);
\node[font=\small] at (2.6,-1.32) {$2$};
\draw[atomLine, line width=0.6pt] (5.2,0) circle (1);
\fill[fmBlue] (5.783,0.212) circle (0.075);
\fill[fmBlue] (4.725,0.399) circle (0.075);
\fill[fmBlue] (5.092,-0.611) circle (0.075);
\node[font=\small] at (5.2,-1.32) {$3$};
\draw[atomLine, line width=0.6pt] (7.800000000000001,0) circle (1);
\fill[fmBlue] (8.383,0.212) circle (0.075);
\fill[fmBlue] (7.588,0.583) circle (0.075);
\fill[fmBlue] (7.217,-0.212) circle (0.075);
\fill[fmBlue] (8.012,-0.583) circle (0.075);
\node[font=\small] at (7.800000000000001,-1.32) {$4$};
\draw[atomLine, line width=0.6pt] (10.4,0) circle (1);
\fill[fmBlue] (10.983,0.212) circle (0.075);
\fill[fmBlue] (10.508,0.611) circle (0.075);
\fill[fmBlue] (9.925,0.399) circle (0.075);
\fill[fmBlue] (9.817,-0.212) circle (0.075);
\fill[fmBlue] (10.292,-0.611) circle (0.075);
\fill[fmBlue] (10.875,-0.399) circle (0.075);
\node[font=\small] at (10.4,-1.32) {$6$};
\draw[atomLine, line width=0.6pt] (0.0,-2.75) circle (1);
\fill[fmBlue] (0.583,-2.538) circle (0.075);
\node[fmRed, font=\small] at (-0.583,-2.962) {$\times$};
\node[font=\small] at (0.0,-4.07) {$\br{1}$};
\draw[atomLine, line width=0.6pt] (2.6,-2.75) circle (1);
\node[fmRed, font=\small] at (3.183,-2.538) {$\times$};
\fill[fmBlue] (3.183,-2.538) circle (0.075);
\node[font=\small] at (2.6,-4.07) {$\br2=m$};
\draw[atomLine, line width=0.6pt] (5.2,-2.75) circle (1);
\fill[fmBlue] (5.783,-2.538) circle (0.075);
\node[fmRed, font=\small] at (5.308,-2.139) {$\times$};
\fill[fmBlue] (4.725,-2.351) circle (0.075);
\node[fmRed, font=\small] at (4.617,-2.962) {$\times$};
\fill[fmBlue] (5.092,-3.361) circle (0.075);
\node[fmRed, font=\small] at (5.675,-3.149) {$\times$};
\node[font=\small] at (5.2,-4.07) {$\br{3}$};
\draw[atomLine, line width=0.6pt] (7.800000000000001,-2.75) circle (1);
\fill[fmBlue] (8.383,-2.538) circle (0.075);
\node[fmRed, font=\small] at (7.588,-2.167) {$\times$};
\fill[fmBlue] (7.217,-2.962) circle (0.075);
\node[fmRed, font=\small] at (8.012,-3.333) {$\times$};
\node[font=\small] at (7.800000000000001,-4.07) {$\br{4}$};
\draw[atomLine, line width=0.6pt] (10.4,-2.75) circle (1);
\node[fmRed, font=\small] at (10.983,-2.538) {$\times$};
\fill[fmBlue] (10.983,-2.538) circle (0.075);
\node[fmRed, font=\small] at (9.925,-2.351) {$\times$};
\fill[fmBlue] (9.925,-2.351) circle (0.075);
\node[fmRed, font=\small] at (10.292,-3.361) {$\times$};
\fill[fmBlue] (10.292,-3.361) circle (0.075);
\node[font=\small] at (10.4,-4.07) {$\br{6}$};

\end{tikzpicture}
\caption{Действие осей симметрии, перпендикулярных плоскости проекции:
верхний ряд~— поворотные оси 1, 2, 3, 4, 6; нижний~— инверсионные оси
$\br1$, $\br2=m$, $\br3$, $\br4$, $\br6$. Точка~— направление, выходящее в
верхнюю полусферу, крестик~— в нижнюю (расчёт)}
\label{fig:axes-proj}
\end{figure}

\begin{figure}[ht]
\centering
\begin{tikzpicture}[cubeproj, scale=2.3]
  \CubeEdges
  \draw[fmRed, very thick] (0.5,0.5,-0.35) -- (0.5,0.5,1.35);
  \node[draw=fmRed, fill=fmRed!25, minimum size=2.4mm, inner sep=0] at (0.5,0.5,1.35) {};
  \node[font=\scriptsize, fmRed, anchor=west] at (0.55,0.5,1.4) {4};
  \draw[fmBlue, very thick] (-0.25,-0.25,-0.25) -- (1.25,1.25,1.25);
  \node[regular polygon, regular polygon sides=3, draw=fmBlue, fill=fmBlue!25, minimum size=3.2mm, inner sep=0] at (1.25,1.25,1.25) {};
  \node[font=\scriptsize, fmBlue, anchor=west] at (1.32,1.25,1.25) {3};
  \draw[fmGreen!60!black, very thick] (-0.25,0.5,-0.25) -- (1.25,0.5,1.25);
  \node[ellipse, draw=fmGreen!60!black, fill=fmGreen!25, minimum width=3.2mm, minimum height=1.8mm, inner sep=0] at (1.25,0.5,1.25) {};
  \node[font=\scriptsize, fmGreen!60!black, anchor=west] at (1.32,0.5,1.25) {2};
\end{tikzpicture}
\caption{Оси симметрии куба: три оси 4 (вдоль $\nfam{100}$, показана одна),
четыре оси 3 (вдоль $\nfam{111}$) и шесть осей 2 (вдоль $\nfam{110}$)}
\label{fig:cube-axes}
\end{figure}

\paragraph{Матрицы поворотов}
Правило из п.~\ref{ssec:orth}: столбцы матрицы~— координаты векторов $[100]$,
$[010]$, $[001]$ после поворота. Их удобно находить по стандартной
проекции (рис.~\ref{fig:cubic-dirs}\,\textit{б}).
\begin{itemize}
\item Поворот на $90^\circ$ вокруг $[100]$ (против часовой стрелки, если
смотреть с конца оси): $[100]\to[100]$, $[010]\to[001]$, $[001]\to[0\br10]$:
\[
C_4^{[100]}=\begin{pmatrix}1&0&0\\0&0&\br1\\0&1&0\end{pmatrix},\quad
\operatorname{tr}=1\ \Rightarrow\ \alpha=90^\circ .
\]
\item Поворот на $120^\circ$ вокруг $[1\br11]$: $[100]\to[001]$,
$[010]\to[\br100]$, $[001]\to[0\br10]$:
\[
C_3^{[1\br11]}=\begin{pmatrix}0&\br1&0\\0&0&\br1\\1&0&0\end{pmatrix},\quad
\operatorname{tr}=0\ \Rightarrow\ \alpha=120^\circ ;
\]
по~\eqref{eq:axis-angle} ось $(0+1):(0-1):(0+1)=[1\br11]$.
\item Поворот на $180^\circ$ вокруг $[011]$: $[100]\perp[011]$ переходит в
$[\br100]$, а $[010]$ и $[001]$ меняются местами:
\[
C_2^{[011]}=\begin{pmatrix}\br1&0&0\\0&0&1\\0&1&0\end{pmatrix},\quad
\operatorname{tr}=-1\ \Rightarrow\ \alpha=180^\circ .
\]
\end{itemize}

\begin{vnimanie}[Ось поворота на $180^\circ$]
Для $\alpha=180^\circ$ матрица симметрична, и формула~\eqref{eq:axis-angle}
даёт $0:0:0$. В этом случае ось находят как любой ненулевой столбец
матрицы $A+I$: для $C_2^{[011]}$ это $(0,1,1)$.
\end{vnimanie}

\subsection{Инверсионные оси}\label{ssec:inv-axes}

Поворот можно совместить с отражением в точке на оси (инверсией).
Получаются \term{инверсионные оси} тех же порядков: $\br1$, $\br2$, $\br3$, $\br4$,
$\br6$. Инверсия в начале координат (центр инверсии)
\begin{equation}\label{eq:inversion}
\bar1=\begin{pmatrix}\br1&0&0\\0&\br1&0\\0&0&\br1\end{pmatrix},\qquad
R_i=\bar1R=\begin{pmatrix}\br1&0&0\\0&-\cos\alpha&\sin\alpha\\0&-\sin\alpha&-\cos\alpha\end{pmatrix},
\end{equation}
$\det R_i=-1$, $\operatorname{tr}R_i=-1-2\cos\alpha$. Разберём оси по очереди
(рис.~\ref{fig:axes-proj}, нижний ряд).
\begin{itemize}
\item $\br1$ ($\alpha=0$)~— просто \emph{центр симметрии} (инверсии).
\item $\br2$ вдоль $[100]$: $C_2\bar1=\operatorname{diag}(\br1,1,1)$~— зеркальное
отражение в плоскости $yz$. Значит, \vazhno{ось $\br2$ эквивалентна
перпендикулярной ей плоскости симметрии} $m$ (от англ.\ mirror).
\item $\br3=3+\br1$: ось 3 и центр симметрии.
\item $\br6=3/m$: ось 3 и перпендикулярная ей плоскость.
\item $\br4$~— самостоятельный элемент; он содержит ось 2 (дважды
применённый $\br4$~— поворот на $180^\circ$).
\end{itemize}
Независимы, таким образом, $\br1$, $\br2=m$ и $\br4$; оси $\br3$ и $\br6$~— составные.

\subsection{Зеркально-поворотные оси. Полярность}\label{ssec:mirror-axes}

Несобственный поворот можно получить и иначе~— сочетанием поворота с
отражением в перпендикулярной плоскости. Это \term{зеркально-поворотные}
оси $\tilde1$, $\tilde2$, $\tilde3$, $\tilde4$, $\tilde6$ с матрицей
$R_m=mR$, $\det R_m=-1$, $\operatorname{tr}R_m=-1+2\cos\alpha$. Новых элементов они не
дают:
\[
\tilde1=\br2=m,\quad\tilde2=\br1,\quad\tilde3=\br6,\quad\tilde4=\br4,\quad\tilde6=\br3 .
\]
Инверсионные оси приняты в международной символике, зеркально-поворотные~—
в символике Шёнфлиса ($S_n$) и Шубникова.

\paragraph{Полярность}
Вдоль прямой можно двигаться в две стороны. Если оба направления
равноценны (их совмещает какая-либо операция симметрии~— плоскость,
перпендикулярная прямой, центр симметрии на ней, ось 2,
перпендикулярная ей), прямая \emph{неполярна}; иначе~— \term{полярна}.
Инверсионные оси всегда неполярны; поворотная ось полярна, если её не
пересекают элементы $\br1$, $m$, $2$. Полярности направлений отвечает
полярность физических свойств (пироэлектричество, пьезоэлектричество,
п.~\ref{ssec:point-groups}).

\subsection{Элементы теории групп}\label{ssec:groups}

Комбинации элементов симметрии удобно изучать с помощью теории групп.

\begin{opredelenie}[Группа]
Множество элементов $G=\{g_1,g_2,\ldots\}$ с операцией «умножения», для
которой:
\begin{enumerate}
\item произведение любых двух элементов~— элемент группы: $g_ig_j=g_k$;
\item умножение ассоциативно: $g_i(g_jg_l)=(g_ig_j)g_l$;
\item есть единственный единичный элемент $e$: $eg_i=g_ie=g_i$;
\item у каждого элемента есть обратный: $g_ig_i^{-1}=e$.
\end{enumerate}
\end{opredelenie}

«Умножение» не обязательно арифметическое: в группах симметрии это
последовательное выполнение операций ($g_ig_j$~— сначала $g_j$, потом $g_i$).
Коммутативность не требуется; группы, в которых $g_ig_j=g_jg_i$,
называют \term{абелевыми}.

\paragraph{Основные понятия}
\begin{itemize}
\item \term{Порядок группы}~— число её элементов (конечное или
бесконечное).
\item \term{Подгруппа}~— часть элементов, сама образующая группу; содержит
$e$. По \emph{теореме Лагранжа} порядок подгруппы~— делитель порядка
группы; поэтому группа простого порядка не имеет нетривиальных подгрупп.
\item \term{Порядок элемента} $g$~— наименьшее $r$, при котором $g^r=e$;
степени $g,g^2,\ldots,g^r=e$ образуют \emph{период} элемента. Если вся группа
исчерпывается степенями одного элемента, она \term{циклическая}, а этот
элемент~— её \emph{генератор}; иначе выделяют несколько генераторов.
\item Группы \term{изоморфны}, если между их элементами есть взаимно
однозначное соответствие, сохраняющее умножение,~— с абстрактной точки
зрения это одна группа с разными «именами» элементов. При соответствии
без взаимной однозначности (одному элементу $H$ отвечает несколько
элементов $G$) группы \term{гомоморфны}.
\item Элемент $g_k$ \term{сопряжён} элементу $g_i$, если $g_k=g_j^{-1}g_ig_j$ для
какого-то $g_j$ из группы. Взаимно сопряжённые элементы образуют
\term{класс}; единица всегда образует класс сама по себе. Геометрически
сопряжённые операции~— «одинаковые» операции, переводимые друг в друга
симметрией группы.
\end{itemize}

\begin{primer*}[Примеры групп]
\begin{enumerate}
\item Целые числа по сложению: единица~— $0$, обратный к $n$~— $-n$;
бесконечная абелева группа.
\item Рациональные числа, \emph{кроме нуля}, по умножению: единица~— $1$,
обратный~— $1/n$; бесконечная абелева группа\isprmark.
\item Невырожденные матрицы $n\times n$ по умножению; бесконечная
неабелева группа.
\item Повороты тела вокруг всевозможных осей: единица~— отсутствие
поворота, обратный~— поворот в обратную сторону; бесконечная неабелева
группа.
\end{enumerate}
\end{primer*}
\isprtext{В оригинале к группе по умножению отнесено всё множество
рациональных чисел $p/q$, $q\ne0$; но у нуля нет обратного элемента, поэтому
его нужно исключить ($p\ne0$).}

\paragraph{Таблица умножения (таблица Кэли)}
Элементы выписывают по вертикали и горизонтали; на пересечении строки
$g_i$ и столбца $g_j$ ставят произведение $g_ig_j$. Абстрактных групп
4-го порядка всего две:
\[
\begin{array}{c|cccc}
 & e & a & b & c\\\hline
e & e & a & b & c\\
a & a & b & c & e\\
b & b & c & e & a\\
c & c & e & a & b
\end{array}
\qquad\qquad
\begin{array}{c|cccc}
 & e & a & b & c\\\hline
e & e & a & b & c\\
a & a & e & c & b\\
b & b & c & e & a\\
c & c & b & a & e
\end{array}
\]
Слева~— \emph{циклическая} группа: $a^2=b$, $a^3=c$, $a^4=e$ (генератор $a$);
её реализация~— повороты вокруг оси 4. Справа~— \emph{четверная группа Клейна}:
$a^2=b^2=c^2=e$, $ab=c$; генераторами могут служить любые два из $a$, $b$,
$c$; её реализация~— группа $222$ (три перпендикулярные оси 2) или группа
$2/m$. Обе группы абелевы, и каждый элемент образует свой класс.
Геометрически разные группы симметрии (у одной элемент~— отражение, у
другой~— поворот на $180^\circ$) могут иметь одинаковую таблицу умножения,
т.\,е. быть изоморфными.
